\documentclass[12pt,a4paper]{article}
\usepackage[utf8]{inputenc}
\usepackage[T2A]{fontenc}
\usepackage[russian]{babel}
\usepackage{amsmath,amssymb}
\usepackage{geometry}
\usepackage{booktabs}
\usepackage{hyperref}
\geometry{margin=2.4cm}

\title{Гипотеза общего спектрального ядра для \texttt{TOE.pdf} и \texttt{UGSM}}
\author{}
\date{}

\begin{document}

\maketitle

\begin{abstract}
В этой записке исследуется гипотеза, что базовые объекты теоретических корпусов \texttt{TOE.pdf} и \texttt{ugsm\_dynamics\_audit-3.pdf} могут быть сведены к одному более общему спектральному объекту в смысле протокола Shadow-to-Theory (S2T). Предлагается минимальный мост: корреляционный оператор TOE интерпретируется как тепловое полугрупповое ядро того же оператора Лаплас/Дирак-типа, спектральные моменты которого лежат в основе UGSM. Далее проводится простой численный эксперимент на \(S^3 \times S^1\), показывающий, что первые коэффициенты теплового следа действительно воспроизводят геометрические масштабы вида \(4\pi^3\) и, после естественной нормировки, \(\pi^2\). Третий линейный по \(\pi\) вклад интерпретируется как кандидат на тополого-голономный или систолический сектор. Вывод носит исследовательский, а не доказательный характер.
\end{abstract}

\section{Вопрос исследования}

Если смотреть на \texttt{TOE.pdf} и UGSM через протокол S2T, возникает естественный вопрос второго уровня: являются ли эти конструкции независимыми теориями или двумя разными проекциями одного более общего спектрального описания? В строгой формулировке вопрос звучит так:
\begin{quote}
существует ли единый базовый оператор, из которого корреляционный объект TOE и спектральный EFT-блок UGSM получаются как разные функциональные представления?
\end{quote}

В этой записке предлагается минимальная рабочая гипотеза, которая делает такой вопрос математически осмысленным.

\section{Выбор базовых постулатов}

\subsection{Постулаты TOE}

Из \texttt{TOE.pdf} берём следующие базовые пункты:
\begin{enumerate}
    \item фундаментальным объектом является компактный самосопряжённый нелокальный корреляционный оператор \(\hat C\);
    \item ядро \(\hat C\) имеет гауссову тепловую форму;
    \item спектральное действие строится как след функции от базового оператора;
    \item геометрия возникает из локальной структуры корреляционного ядра.
\end{enumerate}

Содержательно это означает, что TOE уже почти прямо связывает физику с тепловым полугрупповым оператором типа
\[
\hat C_\sigma \sim e^{-\sigma \mathcal{L}/2},
\]
где \(\mathcal{L}\) --- оператор Лаплас/Дирак-типа.

\subsection{Постулаты UGSM}

Из \texttt{ugsm\_dynamics\_audit-3.pdf} берём следующие базовые пункты:
\begin{enumerate}
    \item минимальный вакуумный сектор задаётся на \(M = S^3 \times S^1\);
    \item центральный объект --- спектральный функционал
    \[
    S_{spec}(\Lambda)=\mathrm{Tr}\, f(D_M^2/\Lambda^2);
    \]
    \item наблюдаемые выражаются через спектральные инварианты и их усечённые комбинации;
    \item по ходу документа вводится единая цепочка вида
    \[
    f \to B_j \to C_2 \to R_2 \to S_{vac}.
    \]
\end{enumerate}

Это означает, что UGSM по существу работает не с иным типом объекта, а с иным \emph{функциональным представлением} спектра того же рода операторов.

\section{Кандидат на математический мост}

Минимальная гипотеза общего ядра выглядит так:
\begin{equation}
\hat C_\sigma = e^{-\sigma D_M^2/2}.
\end{equation}

Тогда спектральный функционал UGSM получается как интегральное преобразование от того же объекта. Если функция \(f\) допускает представление Лапласа
\begin{equation}
 f(x)=\int_0^{\infty} \mu(t)e^{-tx}\,dt,
\end{equation}
то
\begin{equation}
\mathrm{Tr}\, f(D_M^2/\Lambda^2)
= \int_0^{\infty} \mu(t)\, \mathrm{Tr}\, e^{-t D_M^2/\Lambda^2}dt.
\end{equation}

После подстановки \(\sigma = 2t/\Lambda^2\) имеем
\begin{equation}
S_{spec}(\Lambda)
=\int_0^{\infty} \mu(t)\, \mathrm{Tr}\, \hat C_{2t/\Lambda^2}\, dt.
\end{equation}

Именно это равенство и есть минимальный математический мост: 
\begin{quote}
UGSM можно интерпретировать как спектрально-моментную редукцию того же корреляционного объекта, который в TOE считается фундаментальным.
\end{quote}

Если эта гипотеза верна, тогда различие между TOE и UGSM состоит не в базовом объекте, а в уровне описания:
\begin{itemize}
    \item TOE акцентирует фундаментальный нелокальный коррелятор и глобальную архитектуру;
    \item UGSM акцентирует его спектральные моменты, усечения и EFT-интерпретацию на конкретном компактном секторе.
\end{itemize}

\section{Физическая интерпретация моста}

В таком чтении общее физическое содержание выглядит следующим образом:
\begin{enumerate}
    \item TOE задаёт ``полную тень'' через корреляционный оператор.
    \item UGSM выделяет из этой тени компактный спектральный сектор на \(S^3 \times S^1\).
    \item Спектральные коэффициенты Seeley--DeWitt становятся теми самыми геометрическими инвариантами, из которых собираются наблюдаемые комбинации.
    \item Различные численные формулы UGSM тогда перестают быть независимыми ``угаданными'' выражениями и становятся приближёнными проекциями общей спектральной структуры.
\end{enumerate}

Это не доказывает эквивалентность двух теорий, но показывает, что их базовые объекты \emph{могут быть сведены к одному операторному классу}.

\section{Численный эксперимент}

Чтобы проверить, не является ли такой мост чистой риторикой, проведём минимальный спектральный эксперимент на \(S^3 \times S^1\) при радиусах \(R_3 = R_1 = 1\). Берём скалярный лапласиан с собственными значениями
\[
\lambda_{n,m}=n(n+2)+m^2,
\]
где кратность сферической части равна \((n+1)^2\), а \(m\in\mathbb Z\).

Рассматриваем тепловой след
\[
H(t)=\mathrm{Tr}\, e^{-t\Delta},
\]
и аппроксимируем малотемпературную асимптотику в виде
\[
(4\pi t)^2 H(t) \approx a_0 + a_2 t + a_4 t^2.
\]

Численный пересчёт выполнен в файле \texttt{spectral\_bridge\_experiment.py}; результаты сохранены в \texttt{spectral\_bridge\_results.json}.

\subsection{Результаты}

Получено:
\begin{center}
\begin{tabular}{lcc}
\toprule
Величина & Численный fit & Теоретическое значение \\
\midrule
$a_0$ & 124.0268873911 & $4\pi^3 = 124.0251067212$ \\
$a_2$ & 123.8993672201 & $4\pi^3 = 124.0251067212$ \\
$a_2/(4\pi)$ & 9.8595983695 & $\pi^2 = 9.8696044011$ \\
\bottomrule
\end{tabular}
\end{center}

Точность достаточна для содержательного вывода: первые два коэффициента действительно воспроизводят геометрические масштабы вида
\[
4\pi^3, \qquad \pi^2.
\]

Это важно, потому что именно такие масштабы фигурируют в UGSM-подобных структурных комбинациях.

\subsection{Интерпретация линейного вклада \texorpdfstring{$\pi$}{pi}}

Линейный вклад \(\pi\) в UGSM не извлекается из данного простого эксперимента автоматически. Однако есть естественный кандидат: фундаментальная длина цикла \(S^1\), то есть его систола,
\[
\mathrm{sys}(S^1)=2\pi R_1.
\]
При \(R_1=1\) половина этой длины даёт
\[
\frac{\mathrm{sys}(S^1)}{2}=\pi.
\]

Поэтому рабочая гипотеза выглядит так:
\begin{itemize}
    \item \(4\pi^3\) --- объёмный \(a_0\)-сектор;
    \item \(\pi^2\) --- curvature-sector после нормировки \(a_2/(4\pi)\);
    \item \(\pi\) --- тополого-голономный или систолический сектор, связанный с фундаментальным циклом \(S^1\).
\end{itemize}

Это уже не строгий вывод, а исследовательская гипотеза, но она физически согласует структуру UGSM с геометрией компактного сектора, а не с произвольной нумерологией.

\section{Что этот эксперимент реально показывает}

Эксперимент не доказывает, что TOE и UGSM --- одна и та же теория. Но он показывает три более скромных и важных факта.

\begin{enumerate}
    \item Базовые объекты двух теорий действительно можно поместить в один операторный класс.
    \item Первые геометрические коэффициенты на \(S^3 \times S^1\) естественно воспроизводят именно те масштабы, которые UGSM использует как структурное ядро.
    \item Следовательно, гипотеза ``общей тени'' перестаёт быть чисто метафорической и приобретает конкретную математическую форму.
\end{enumerate}

\section{Что пока не доказано}

Чтобы перейти от гипотезы общего ядра к строгому утверждению, необходимо ещё доказать как минимум следующее:
\begin{enumerate}
    \item что оператор TOE действительно редуцируется к \(e^{-\sigma D_M^2/2}\) на компактном секторе UGSM;
    \item что выбор функции \(f\) в UGSM соответствует корректному преобразованию меры \(\mu(t)\) в TOE;
    \item что линейный \(\pi\)-вклад выводится не эвристически, а из голономии, систолы или Wilson-loop сектора;
    \item что остальные наблюдаемые UGSM можно получить из той же редукции без внешних ручек.
\end{enumerate}

Именно эти пункты отделяют нынешний мост-гипотезу от доказанной редукции одной теории к другой.

\section{Итоговый вывод}

В результате проведённого исследования можно сформулировать осторожный, но содержательный итог.

\begin{quote}
На текущем уровне анализа базовые объекты TOE и UGSM \emph{могут быть сведены} к одному более общему спектральному описанию: корреляционный оператор TOE естественно интерпретируется как тепловое полугрупповое ядро того же оператора Лаплас/Дирак-типа, спектральные моменты которого используются в UGSM. Численный эксперимент на \(S^3 \times S^1\) подтверждает, что первые геометрические коэффициенты действительно производят масштабы \(4\pi^3\) и \(\pi^2\), а линейный \(\pi\)-вклад допускает естественную тополого-геометрическую интерпретацию.
\end{quote}

Следовательно, гипотеза о том, что обе теории являются разными тенями одного более общего спектрального закона, является \emph{математически правдоподобной}. Однако её статус пока следует считать не доказанным результатом, а сильной исследовательской программой для следующего S2T-цикла.

\end{document}